\lecture{9.3}{Causal Dimension Reduction}
\label{lec:9-3}

\subsection{Lower-Dimensional Treatment Representations}

When $A\in\R^p$ is high-dimensional, direct estimation of
$a\mapsto\E\{\Ya{a}\}$ is often infeasible. A natural strategy is to replace
$A$ by a lower-dimensional representation
\[
  G=g(A)\in\R^d,\qquad d<p.
\]
The goal is not merely to compress $A$, but to preserve the causal information
in $A$ that is relevant for $Y$.

\subsection{Effect-Preserving Maps $g(A)$}

\begin{definition}[Effect-preserving representation]
A map $g:\mathcal{A}\rightarrow\mathcal{G}$ is effect-preserving for the mean
counterfactual outcome if
\[
  \E\{\Ya{a}\}=\theta\{g(a)\}
\]
for some function $\theta$ and for all exposure values $a$ in the target set.
\end{definition}

This means that any two exposure values with the same reduced representation
have the same mean counterfactual outcome:
\[
  g(a)=g(a')\quad\Rightarrow\quad \E\{\Ya{a}\}=\E\{\Ya{a'}\}.
\]
The reduced feature $g(A)$ is then sufficient for the causal effect of $A$ on
the mean outcome.

\subsection{Linear Dimension Reduction $g(A;\beta)$}

A common model is a linear reduction:
\[
  g(A;\beta)=\beta^\top A,
\]
or, for $d>1$,
\[
  g(A;\beta)=\beta^\top A,\qquad \beta\in\R^{p\times d}.
\]
The causal dimension-reduction condition becomes
\[
  \E\{\Ya{a}\}=\theta(\beta^\top a).
\]
Here $\beta$ defines directions in the exposure space that are causally relevant
for the outcome.

\begin{figure}[H]
\centering
\begin{tikzpicture}[
  box/.style={draw, rounded corners=2pt, minimum width=2.8cm, minimum height=.8cm},
  arr/.style={-{Latex[length=2.3mm]}, thick},
  node distance=1.1cm
]
\node[box] (A) {High-dimensional\\ exposure $A\in\R^p$};
\node[box, right=of A] (G) {Reduced exposure\\ $g(A)\in\R^d$};
\node[box, right=of G] (Y) {Causal response\\ $\E\{\Ya{a}\}$};
\draw[arr] (A) -- node[above] {$g$} (G);
\draw[arr] (G) -- node[above] {$\theta$} (Y);
\end{tikzpicture}
\caption{Causal dimension reduction seeks a low-dimensional exposure summary that preserves counterfactual means.}
\label{fig:ch9-causal-dimred}
\end{figure}

\subsection{Why PCA Is Insufficient}

Principal component analysis finds directions that explain variation in the
exposure distribution. It does not use the outcome and does not distinguish
causal signal from confounding. A principal component can explain large
variation in $A$ while having no causal effect on $Y$, and a low-variance
component can be highly relevant to $Y$.

More importantly, preserving the association between $A$ and $Y$ is not enough.
Because
\[
  \E\{\Ya{a}\}\ne \E(Y\given A=a)
\]
in the presence of confounding, a dimension-reduction method that optimizes
prediction of $Y$ from $A$ may preserve confounded associations rather than
causal effects.

\subsection{Confounding-Aware Dimension Reduction}

Under conditional exchangeability,
\[
  \E\{\Ya{a}\}=\E\{\E(Y\given A=a,W)\}.
\]
A confounding-aware reduction must therefore preserve the standardized outcome
regression, not merely the raw conditional mean $\E(Y\given A=a)$. The
covariates $W$ enter because the causal response at $a$ is obtained by averaging
over the target population distribution of $W$.

One conceptual approach is to search for $\beta$ such that
\[
  \E\{\E(Y\given A=a,W)\}
  =
  \theta(\beta^\top a)
\]
over the exposure values of interest.

\subsection{Causal Sufficient Dimension Reduction}

\begin{definition}[Causal sufficient dimension reduction]
A reduction $g(A)$ is causally sufficient for the mean effect of $A$ if the
counterfactual mean depends on $a$ only through $g(a)$:
\[
  \E\{\Ya{a}\}=\E\{\Ya{a'}\}
  \quad\text{whenever}\quad g(a)=g(a').
\]
\end{definition}

This definition is about counterfactual distributions, not the observed exposure
distribution. Estimating such a reduction requires causal assumptions, adequate
support, and a method that accounts for confounding.

\begin{keyidea}
Causal dimension reduction is not ordinary compression. It asks for a
low-dimensional representation that preserves intervention effects, not just
variation in the exposure or prediction of the observed outcome.
\end{keyidea}
